% © 2026 Ing. David Jaroš — CC BY-NC-ND 4.0
%
% This work is licensed under a Creative Commons Attribution-NonCommercial-NoDerivatives
% 4.0 International License (CC BY-NC-ND 4.0).
%
% B_base = N_eff^{3/2} — Approach A1: Spinor Degrees of Freedom
% Track: CORE — α derivation — v56 baseline
% Date: 2026-03-07

\ifdefined\INCLUDEMODE
  % Being included in another document — skip preamble
\else
  \documentclass[11pt,a4paper]{article}
  \usepackage{amsmath,amssymb,amsthm}
  \usepackage{hyperref}

  \newtheorem{theorem}{Theorem}
  \newtheorem{proposition}{Proposition}
  \newtheorem{conjecture}{Conjecture}
  \theoremstyle{remark}
  \newtheorem{remark}{Remark}

  \title{Approach A1: $B_{\mathrm{base}} = N_{\mathrm{eff}}^{3/2}$\\
    \large Spinor Degrees of Freedom in $\mathbb{C}\otimes\mathbb{H}$}
  \author{Ing.\ David Jaroš}
  \date{2026-03-07}

  \begin{document}
  \maketitle
\fi

\section{Approach A1: Spinor Degrees of Freedom}
\label{sec:spinor_approach}

\textbf{Status: [MOTIVATED CONJECTURE]} — algebraically self-consistent;
not a rigorous proof. No circular use of $\alpha^{-1} = 137$.

\subsection{Starting Point}

The proved one-loop result is
\begin{equation}
  B_0 = \frac{2\pi N_{\mathrm{eff}}}{3} \approx 25.13,
  \qquad N_{\mathrm{eff}} = 12,
  \label{eq:B0}
\end{equation}
while the phenomenologically required value is $B_{\mathrm{req}} \approx 46.3$,
suggesting a gap factor $B_{\mathrm{req}} / B_0 \approx 1.84$.
The hypothesis here is that $B_{\mathrm{base}} = N_{\mathrm{eff}}^{3/2}$ arises
naturally from the spinor structure of the biquaternionic field.

\subsection{Algebraic Setup}

The biquaternionic field $\Theta$ takes values in $\mathbb{C}\otimes\mathbb{H}
\cong \mathrm{Mat}(2,\mathbb{C})$.  A single mode $\Theta_k$ is therefore
a $2\times 2$ complex matrix — i.e.\ it carries $4$ complex, or
$8$ real, components.  Under the decomposition
\begin{equation}
  \Theta_k = \Theta_k^R + i\,\Theta_k^I,
  \qquad \Theta_k^{R,I} \in \mathbb{H},
\end{equation}
each quaternionic part $\Theta_k^{R,I}$ has $4$ real components.

The effective mode count established in UBT is $N_{\mathrm{eff}} = 12$, arising
from $N_{\mathrm{phases}} \times N_{\mathrm{helicity}} \times N_{\mathrm{charge}}
= 3 \times 2 \times 2$.  Each mode $\Theta_k$ thus contributes one unit to
$N_{\mathrm{eff}}$.

\subsection{Vacuum Polarisation Tensor}

The vacuum polarisation tensor $\Pi_{\mu\nu}$ is built from bilinear products of
$\Theta_k$.  In the biquaternionic sector its entries are $\mathrm{Mat}(2,\mathbb{C})$
valued, giving a naive count of $N_{\mathrm{eff}}^2 = 144$ independent components.

However, the Ward identity $q^\mu \Pi_{\mu\nu}(q) = 0$ and the Hermiticity
condition $\Pi_{\mu\nu}^\dagger = \Pi_{\nu\mu}$ impose constraints.  In
$d=4$ dimensions these reduce the independent transverse degrees of freedom
by a factor proportional to the square root of the field-space dimension.

\textbf{Heuristic counting:}
\begin{itemize}
  \item Mode $\Theta_k$ lives in $\mathrm{Mat}(2,\mathbb{C})$ — a space of
    (complex) dimension $4 = N_{\mathrm{eff}}/3$ per mode.
  \item For $N_{\mathrm{eff}}$ modes the vacuum polarisation has
    $N_{\mathrm{eff}}^2$ naive entries.
  \item Unitarity of the $S$-matrix imposes the optical theorem on $\Pi_{\mu\nu}$,
    effectively projecting onto a subspace.  In a theory where the field
    transforms in the adjoint of a group of rank $r$, the surviving independent
    components scale as $\sqrt{r}$ times the dimension.
  \item Here the relevant ``rank'' is provided by the $N_{\mathrm{eff}}$
    independent spinorial modes: the unitarity constraint selects
    $\sqrt{N_{\mathrm{eff}}}$ combinations per mode.
\end{itemize}

Combining:
\begin{equation}
  B_{\mathrm{base}} \sim N_{\mathrm{eff}} \times \sqrt{N_{\mathrm{eff}}}
    = N_{\mathrm{eff}}^{3/2}.
  \label{eq:B_spinor_heuristic}
\end{equation}
For $N_{\mathrm{eff}} = 12$:
\begin{equation}
  B_{\mathrm{base}} = 12^{3/2} = \sqrt{1728} \approx 41.57.
\end{equation}

\subsection{QED Consistency Check}

For a single charged Dirac fermion ($N_{\mathrm{eff}} = 1$):
\begin{equation}
  B_{\mathrm{base}}\big|_{N=1} = 1^{3/2} = 1.
\end{equation}
The one-loop vacuum polarisation coefficient in QED (Feynman gauge) is
$\beta_1 = 1/3$, corresponding to $B_0 = 2\pi/3 \approx 2.09$.
The spinor heuristic gives $B_{\mathrm{base}} = 1$ — a different quantity.
These are not identical, but $B_{\mathrm{base}}$ here is the
\emph{mode-counting prefactor} before the $2\pi/3$ normalisation, so there
is no direct contradiction; rather $B = B_0 \times (B_{\mathrm{base}}/N_{\mathrm{eff}})$
would need to be verified with $N_{\mathrm{eff}} = 1$.

\textbf{Gap not closed by this argument alone.}
The argument replaces ``$N_{\mathrm{eff}}^{3/2}$'' as a mysterious input with a
heuristic counting claim about unitarity projections. It is not a derivation
from the UBT Lagrangian — it identifies a plausible algebraic \emph{origin}
for the exponent $3/2$.

\subsection{Conclusion}

\begin{conjecture}[Spinor counting — MOTIVATED CONJECTURE]
  In the biquaternionic vacuum polarisation, unitarity constraints on
  $\mathrm{Mat}(2,\mathbb{C})$-valued bilinears reduce $N_{\mathrm{eff}}^2$
  naive components to $N_{\mathrm{eff}}^{3/2}$ independent transverse degrees
  of freedom, yielding
  \begin{equation}
    B_{\mathrm{base}} = N_{\mathrm{eff}}^{3/2} = 12^{3/2} \approx 41.57.
    \label{eq:B_conjecture_spinor}
  \end{equation}
\end{conjecture}

\begin{remark}
  To promote this to a proof one must:
  \begin{enumerate}
    \item Compute $\Pi_{\mu\nu}$ explicitly from $\mathcal{L}[\Theta]$ in
      the biquaternionic sector.
    \item Count independent transverse components after Ward identities.
    \item Show the surviving count is exactly $N_{\mathrm{eff}}^{3/2}$, not
      an approximation.
  \end{enumerate}
  Until steps 1--3 are completed this remains \textbf{[MOTIVATED CONJECTURE]},
  not \textbf{[PROVED]}.
\end{remark}

\ifdefined\INCLUDEMODE\else
\end{document}
\fi
